\documentclass{article}

\usepackage[a4paper]{geometry}
\usepackage[english]{babel}
\usepackage{parskip}

\usepackage{amsmath, amssymb}
\usepackage{graphicx}
\usepackage{subcaption}
\usepackage{float}
\usepackage[hidelinks]{hyperref}

\usepackage{listings}
\usepackage{xcolor}

\definecolor{codegreen}{rgb}{0,0.6,0}
\definecolor{codegray}{rgb}{0.5,0.5,0.5}
\definecolor{codepurple}{rgb}{0.58,0,0.82}
\definecolor{backcolour}{rgb}{0.95,0.95,0.92}

\lstdefinestyle{mystyle}{
  backgroundcolor=\color{backcolour},
  commentstyle=\color{codegreen},
  keywordstyle=\color{magenta},
  numberstyle=\tiny\color{codegray},
  stringstyle=\color{codepurple},
  basicstyle=\ttfamily\footnotesize,
  breaklines=true,
  captionpos=b,
  keepspaces=true,
  numbers=left,
  numbersep=5pt,
  showstringspaces=false,
  tabsize=2
}
\lstset{style=mystyle}

\newcommand{\CodeDir}{Code}
\newcommand{\PlotDir}{Plots}
\newcommand{\CalcDir}{Calculations}

\title{NURa Hand-in 1}
\author{Name (Student number)}
\date{Date}

\begin{document}
\maketitle

\textit{Write a short description of your solution in each section. Explain your choices. 
Explicitly draw conclusions from your results.}

\textit{You can refer to selected lines of code where you think it is necessary. 
However, you must provide the full code relevant to each (sub)question below it.}

\textit{Don't forget to describe your figures and 
mention what conclusions you draw from them!}


\section{Satellite galaxies around a massive central}

\subsection{(1a) Determining the normalization}
We will first solve the integral 
\begin{equation}
\int \int \int_V n(x) dV.
\end{equation}
Afterwards we will solve for A by requiring that the integral equals $\langle N_\mathrm{sat} \rangle$.
Note that
\begin{equation}
n(x) = A \langle N_\mathrm{sat} \rangle 
\left( \frac{x}{b}\right)^{a-3}
\exp\left[-\left(\frac{x}{b}\right)^c\right]
\end{equation}
does not depend on $\theta$ and  $\phi$. Using spherical coordinates turns this 3d integral into a 
1-dimensional integral which is a lot easier to solve numerically. 
The integral then becomes
\begin{equation}
4\pi b^2\int_{10^{-5}}^5 
A \langle N_\mathrm{sat} \rangle 
\left( \frac{x}{b}\right)^{a-1}
\exp\left[-\left(\frac{x}{b}\right)^c\right]
dx.
\end{equation}
We will use Romberg integration as our method of numeric integration.
Romberg integration is a lot more accurate if the function does not vary rapidly and is close to linear,
as it uses evenly spaced samples and trapezoid to approximate the function. 
We rewrite the integral into an integral over logspace, as this makes sure the inner region of the profile, 
where the function varies the most and the region which contributes the most, gets a sufficient amount of samples.
So we make the substitution $u \equiv \mathrm{ln}x$, $\exp(u)du \equiv dx$
and rewrite the integral to:
\begin{equation}
4\pi b^{3 - a}\int_{\ln(10^{-5})}^{\ln(5)} 
A \langle N_\mathrm{sat} \rangle 
\exp\left(au
-\frac{\exp(cu)}{b^c}\right)
du.
\end{equation}

We solve the integral first by using $A=1$. We can then rescale the solution so that it equals $\langle N_\mathrm{sat} \rangle$
since it is linear in A. In other words, solve the integral for $A=1$, then $A = N_\mathrm{sat} / \mathrm{integral}$.

We find: $A = $
\input{Calculations/satellite_A.txt}

\subsubsection{Code for integrating}
\lstinputlisting[language=Python]{\CodeDir/integration.txt}

\subsection{(1b) Sampling the satellite distribution}

Describe the sampling method used (e.g. rejection sampling).  
Explain how the random number generator was implemented and how its randomness was tested.

To go from $n(x)$ to $N(x)$ we can use that $n(x)dV = 4\pi x^2 n(x) dx$.
But note that this is exactly the \textit{number} of galaxies in the infinitesimal range $[x, x+dx)$.
So $N(x) = 4\pi x^2 n(x)$ from which we find that
\begin{equation}
  p(x)dx = \frac{4\pi x^2}{ \langle N_\mathrm{sat}\rangle } n(x) dx.
\end{equation}
Filling in $n(x)$:
\begin{equation}
  p(x)dx = 
  4\pi b^2
  A 
  \left( \frac{x}{b}\right)^{a-1}
  \exp\left[-\left(\frac{x}{b}\right)^c\right] 
  dx.
\end{equation}

\begin{figure}[H]
\centering
\includegraphics[width=0.85\textwidth]{\PlotDir/my_solution_1b.png}
\caption{Comparison between the analytical satellite distribution \(N(x)\) and the histogram of 10 000 sampled satellite galaxies in logarithmic bins.}
\label{fig: probability distribution satellites}
\end{figure}

Does your sampled distribution agree with the analytical prediction? How much samples are thrown away because of rejection sampling?

\subsubsection{Code for random number generation}
\lstinputlisting[language=Python]{\CodeDir/rng.txt}

\subsubsection{Code for sampling}
\lstinputlisting[language=Python]{\CodeDir/distribution.txt}
 

\subsection{(1c) Selecting and sorting satellites}

Describe the selection of 100 galaxies (equal probability, no duplicates, no rejection).
Explain the sorting algorithm you used.

I select 100 galaxies from the N samples obtained in question 1b
by employing the Fisher-Yates Shuffle algorithm (see wikipedia.org/Fisher_Yates_shuffle).
This algorithm works as follows:
\begin{itemize}
\item set N equal to the length of the array
\item randomly draw an index: $i \in [0, N-1]$
\item store the element belonging to this index (the chosen element)
\item swap the Nth element of the array with the chosen element
\item Decrease N by 1
\item repeat until desired number of elements are stored
\end{itemize}

Step 2 ensures that each galaxy has equal probability of being chosen.
Step 2, 4 and 5 ensure that there are no duplicates: If a galaxy is drawn,
then it will be put in the Nth slot of the array, call it M. In the next iterations
$N < M$ (due to step 5), so then in step 2 we can never have $i = M-1$, so the galaxy can 
never be selected again.
At no point during this algorithm do we reject a drawn index. 
So this algorithm satisfies the requirements equal probability, no duplicates, and no rejection.

We proceed to sort the selected galaxies based on the distance from the central.
For this we use the quicksort algorithm. We do not need a stable algorithm. 
This sorting algorithm is implemented inside the Sorter class.
It is divided up into two big blocks of code. The first block sorts an array \textbf{without} bookkeeping an index array 
on the side as well.
The second block does to the bookkeeping and then returns and index array as well.
I have separated this into two big blocks instead of doing the if statements inside the for loops because 
I dont want the number of if statement checks that depend on the length of the array. This safes us a lot of operations
especially for larger arrays.

First the algorithm checks if the array is length 1. In this case it is trivially sorted.
Then the algorithm checks if the array is length 2 and sorts it if necessary.
For larger arrays, the quicksort part comes into play.
We choose a pivot index halfway inside the array $(i_\mathrm{pivot} = N//2)$.
We sort the first, pivot, and last element of the array. Which requires 3 if statements at least.
Redefine the pivot element to be the element that is currently located at the pivot index.

We now initiate a for loop which will increase an index i from 0 to the pivot index while also 
decreasing an index j from N-1 to the pivot index (here N is the length of the array).
We keep track whether the ith (and jth) element are greater (smaller) than the pivot element. 
If i (or j) satisfies this condition, we stop increasing i (or stop decreasing j).
if both i and j are locked, we swap the ith and jth element and unlock both i and j.
We also include an if statement to make sure that i (j) can never exceed (subceed) the pivot index.
We continue until i=j. Then all the elements left of the pivot are smaller than the pivot, and all the elements to 
the right of the pivot are greater than the pivot.
We now call the quicksort algorithm on the two subarrays left and right from the pivot. 
When keeping track of an indexing array, we simply do all the swaps on an arange array of length N as well.

In Figure \ref{fig: cumulative number of satellites} we plot the cumulative number of the 100 selected samples
as a function of radius. We see that the steepest increase in the plot is right where the probability distribution 
in \ref{fig: probability distribution satellites} attains its maximum. This is indeed what we expect.
Due to the small number of samples (100) we do not get any information on the cumulative distribution
in the regimes where the probability distribution is small.

\begin{figure}[H]
\centering
\includegraphics[width=0.8\textwidth]{\PlotDir/my_solution_1c.png}
\caption{Cumulative number of selected galaxies as a function of relative radius.}
\label{fig: cumulative number of satellites}
\end{figure}

\subsubsection{Code for selection}
\lstinputlisting[language=Python]{\CodeDir/selection.txt}

\subsubsection{Code for sorting}
\lstinputlisting[language=Python]{\CodeDir/sorter.txt}


\subsection{(1d) Numerical derivative}

The analytic expression for the derivative of the number density as function of scaled radius $x$ is given by
\begin{equation}
\frac{dn}{dx} = 
\frac{A\langle N_\mathrm{sat} \rangle}{b}
\left[
  (a-3)\left(\frac{x}{b}\right)^{a-4} 
 - c\left(\frac{x}{b}\right)^{a+c-4}
\right]
\exp[-\left(\frac{x}{b}\right)^c].
\end{equation}

\begin{table}[H]
\centering
\caption{Derivative of the satellite number density at \(x=1\).}
\begin{tabular}{|c|c|}
\hline
Method & Value \\
\hline
Analytical & \input{\CalcDir/satellite_deriv_analytic.txt} \\
Numerical & \input{\CalcDir/satellite_deriv_numeric.txt} \\
\hline
\end{tabular}
\end{table}

Do your analytical and numerical derivatives agree? Why?

\subsubsection{Code for derivative calculation}
\lstinputlisting[language=Python]{\CodeDir/differentiation.txt}

\subsection{Execution code for exercise 1}
\lstinputlisting[language=Python]{\CodeDir/Q1_SatelliteGalaxy.txt}


\section{Heating and cooling in HII regions}


\subsection{(2a) Photoionization heating and radiative recombination cooling}

Describe the implementation of your root-finding method (e.g. bisection, Newton-Raphson).

We have a function of the form:
\begin{equation}
  f(T) = A - BT + CT\ln (DT),
\end{equation}
where A, B, C, D are independent of T. We are looking for $T \in [1, 10^7]\mathrm{K}$ such that $f(T) = 0$.
I test a few root-finding algorithms: Bisection, false position, and Newton-Raphson. 
For each algorithm I will test the number of iterations required, the time it took to run, and compare their results.

Bisection iteratively splits up a bracket in half and checks which of the two subintervals contain the root. 
It is super consistent, but usually not the fastest. 
It does not take any information from the slope or values of the functions that it finds along the way. 
False position does take into account the values of the function. It uses these to linearly approximate the function
and from this approximation deduce where the root would be. If the estimated root falls outside the initial bracket,
it resorts to bisection instead. 
Lastly, we have Newton-Raphson. An exremely unstable, but potentially really fast algorithm which uses 
the analytic expression for the derivative of the function in question and uses a starting point rather than a bracket. 
In this scenario it is so unstable that it does not converge for most starting values.
Newton-Raphson requires the function to be sufficiently close to linear. 
Also if the second derivative and first derivative have equal sign for a point, this usually means
that Newton-Raphson will fail as it overshoots. 
In order for this method to be effective, the starting point needs to be \textit{sufficiently} close to the true root.
That is why I instead use an improved Newton-Raphson method: I approach the root using false position,
then I switch to Newton-Raphson if the function varies sufficiently in a small enough bracket.


\begin{table}[H]
\centering
\caption{Equilibrium temperature for the simplified heating–cooling model.}
\begin{tabular}{|c|c|c|}
\hline
Temperature [K] & Absolute error & Relative error & Number of iterations & time to complete \\
\hline
\input{\CalcDir/equilibrium_temp_simple_bisection.txt}
\input{\CalcDir/equilibrium_temp_simple_false_position.txt}
\input{\CalcDir/equilibrium_temp_simple_newton_raphson.txt}
\\
\hline
\end{tabular}
\end{table}

How many iterations it took? Discuss the convergence of your root finder.

\subsection{(2b) More realistic model}


\begin{table}[H]
\centering
\caption{Equilibrium temperature for different gas densities.}
\begin{tabular}{|c|c|}
\hline
Density \(n_H\) [cm\(^{-3}\)] & Equilibrium temperature [K] \\
\hline
$10^{-4}$ & \input{\CalcDir/equilibrium_low_density.txt} \\
$1$ & \input{\CalcDir/equilibrium_mid_density.txt} \\
$10^4$ & \input{\CalcDir/equilibrium_high_density.txt} \\
\hline
\end{tabular}
\end{table}

Discuss the results.  What is the difference for the low, intermediate and high gas density?


\subsection{Code for Exercise 2}

\lstinputlisting[language=Python]{\CodeDir/heatingcooling_code.txt} 

\section{Conclusions}

\end{document}
